Analog filters serve as foundational tools in signal processing. They are mathematically constructed to isolate specific frequency components by passing desired signals while attenuating out-of-band noise. When a system requires the extraction of a strict frequency range, a common requirement when isolating faint target signals from overwhelming ambient noise, a bandpass configuration is typically employed. Filter performance relies heavily on the behavior within the passband, where signal loss should ideally be minimal, and the stopbands, which must provide sufficient suppression of unwanted frequencies \cite{riskin2003}. This laboratory exercise focuses on the design, synthesis, and comparative evaluation of four classical analog bandpass approximations: the Butterworth, Chebyshev Type I, Chebyshev Type II, and Elliptic models. The primary goal is to determine how well these architectures can meet a stringent set of specifications. These include passband edges defined at 2000 and 4000~rad/sec with an allowable ripple of 1~dB, alongside stopband edges at 1500 and 4500~rad/sec that demand at least 60~dB of attenuation. To establish a comparative baseline, the physical filter models are evaluated against a mathematically ideal, though physically unrealizable, bandpass response.

Executing this design relies on a systematic computational approach implemented in MATLAB. The initial theoretical step requires calculating the minimum filter order and corresponding natural frequencies necessary to satisfy the required attenuation limits. This is accomplished using specific order-estimation functions (\texttt{buttord}, \texttt{cheb1ord}, \texttt{cheb2ord}, and \texttt{ellipord}). Once the optimal order is established, the transfer functions are synthesized in the continuous $s$-domain by generating their respective numerator and denominator coefficients. Evaluating the frequency response across the target spectrum then involves computing the complex response via the \texttt{freqs} function. From there, magnitude and phase data are extracted and converted into decibels and degrees for standard analysis. As a final analytical step, an ideal bandpass filter is constructed using symbolic variables and the Heaviside step function. This theoretical model acts as a visual and mathematical reference point to better highlight the inherent limitations of the synthesized approximations.